% !PRINTED none

\begin{center}
\gesperrt{Fifth Section.}

\gesperrt{Linear Transformation.}

\rule{0.3\linewidth}{0.4pt}

\S.\,55.

\gesperrt{Introduction of the linear transformation.}
\end{center}

If, in any functions that depend on the $m$ variables $x_1, x_2 \ldots x_m$, we substitute for these variables linear expressions that depend on $m$ new variables $x_1', x_2' \ldots x_m'$
\weq{1}{\begin{aligned}
x_1 &= \alpha_1^{(1)} x_1' + \alpha_1^{(2)} x_2' + \cdots + \alpha_1^{(m)} x_m', \\
x_2 &= \alpha_2^{(1)} x_1' + \alpha_2^{(2)} x_2' + \cdots + \alpha_2^{(m)} x_m', \\
& \cdots\cdots\cdots\cdots\cdots\cdots\cdots\cdots \\
x_m &= \alpha_m^{(1)} x_1' + \alpha_m^{(2)} x_2' + \cdots + \alpha_m^{(m)} x_m',
\end{aligned}}
then we call this a \gesperrt{linear substitution}, and the result a \gesperrt{linear transformation}.

The quantity
\weq{2}{r = \begin{vmatrix}
\alpha_1^{(1)}, & \alpha_1^{(2)} & \ldots & \alpha_1^{(m)} \\
\alpha_2^{(1)}, & \alpha_2^{(2)} & \ldots & \alpha_2^{(m)} \\
\multicolumn{4}{c}{\cdot\ \cdot\ \cdot\ \cdot\ \cdot} \\
\alpha_m^{(1)}, & \alpha_m^{(2)} & \ldots & \alpha_m^{(m)}
\end{vmatrix}}
is called the \gesperrt{substitution determinant}. It must always be different from zero, since otherwise a dependence among the variables $x_1, x_2 \ldots x_m$ would be expressed by (1).

By a linear substitution every homogeneous function of the variables $x$ passes over into a homogeneous function of the same degree in the variables $x'$.
